(1.7) that @@M0@@ is determined by its value on the element @@M1@@ of @@M2@@ (resp.\ \ldots)..\ Suppose that @@M3@@ is a functor such as @@M4@@, etc..; the element @@M5@@ is a formal series with coefficients in @@M6@@; the @@M7@@-th coefficient of this series will be called the \uline{@@M8@@-th universal polynomial of} @@M9@@; knowing these polynomials is therefore equivalent to knowing @@M10@@. By the proof of 1.6, these polynomials are themselves determined by the element @@M11@@, the calculation following the method of the proof: we calculate the elements @@M12@@ by decomposing @@M13@@ into @@M14@@ and using the additivity of @@M15@@; the polynomials thus obtained are symmetric in the @@M16@@ (@@M17@@ fixed and @@M18@@ varying) and can be expressed in terms of the @@M19@@; by gluing with the aid of the filtration of @@M20@@, we thus determine the universal polynomials of @@M21@@. If @@M22@@, the calculation of @@M23@@ is performed by substituting the coefficients @@M24@@ of the @@M25@@ for the @@M26@@ in the universal polynomials.

Suppose there is an integer @@M27@@ such that the coefficient of @@M28@@ in @@M29@@ is a polynomial whose isobaric components have weights ranging between @@M30@@ and @@M31@@; then the @@M32@@-th universal polynomial of @@M33@@ has isobaric components whose weights range between @@M34@@ and @@M35@@: this follows immediately from the calculation described above and from the theorem on symmetric functions.

